\subsection{What current Gibbs-state results do and do not establish}
\label{threshold_physics:literature}

The following sources were checked against primary records on
14 November 2026. They guide proposed sufficient conditions; none is
being cited as a theorem about unknown-input AA simulation.

\begin{itemize}
\item
\href{https://arxiv.org/abs/2403.16850v2}{Bakshi--Liu--Moitra--Tang,
\emph{High-Temperature Gibbs States are Unentangled and Efficiently
Preparable}, v2 (2025)} prove high-temperature separability and classical
preparation for local Hamiltonians. Their updated abstract gives
separability at inverse temperature \(O(1/\mathfrak d)\), and a
preparation regime \(O(1/\mathfrak d^2)\), under the paper's normalized
interaction-degree convention. The older v1/FOCS formulation has
different quantitative parameters; they should not be mixed.

\item
\href{https://arxiv.org/abs/2607.28536v2}{Putterman--Zlokapa--Cotler,
\emph{When quantum thermal states look classical}, v2 (2026)} use
locality \(k\) and per-site Pauli strength \(s\), allowing long-range
interactions. Their abstract and Sections 1.1--1.4 give a universal
separability/classical-preparation scale \(\beta=O(1/(sk))\) and
distinguish it from thermodynamic analyticity and other transitions.
Their measurement-outcome sampling task produces a new random input;
it does not estimate an arbitrary \(Q(x)\) from a few queries to the
particular unknown input \(x\).

\item
\href{https://arxiv.org/abs/2411.03413}{Chen--Chen--Yin--Zhang,
\emph{Rapid Mixing at the Uniqueness Threshold} (2024)} obtain polynomial
Glauber mixing at critical uniqueness thresholds and investigate spectral
independence there. This reinforces the need to state which kind of
threshold is meant. A sampler using polynomially many operations in
\(n\) is compatible with a classical query lower bound growing with
\(n\).

\item
\href{https://journals.aps.org/prl/abstract/10.1103/PhysRevLett.134.199901}
{Kuwahara--Kato--Brand\~ao's 2025 erratum} identifies failures in the
original 2020 effective-Hamiltonian/CMI cluster-expansion proof. Its
recovered and unresolved claims have different subsystem qualifications.
We do not import the original unrestricted theorem.
\href{https://arxiv.org/abs/2504.02235v2}{Kato--Kuwahara, v2 (2025)}
instead develop quantum belief-propagation channels and give CMI decay
under specified locality/mixing or clustering assumptions, with
high-temperature consequences.
\href{https://arxiv.org/abs/2504.02208}{Chen--Rouz\'e (2025)} prove a
local recovery statement at arbitrary temperature for bounded interaction
degree, including CMI decay for a shielded small region. Neither result
is automatically a classical adaptive-query theorem.
\end{itemize}

There are three independent logical gaps to avoid. First, separability
does not suffice: every classical input law is the diagonal of the
separable state \(\sum_x\mu(x)|x\rangle\langle x|\), including the
character-code full-AA counterexample. Second, even easy global
sampling does not reveal the hidden data of the particular oracle
input. The character-code mixture is easy to generate by choosing its
branch and then choosing either uniform bits or a uniform hidden code
label, yet it has a one-versus-\(\Omega(\log n)\) query separation.
Third, small quantum CMI is not automatically preserved by measuring
the conditioning system: information processing on the middle register
can increase conditional mutual information. The classical measured
conditional laws must be controlled in their own right.

In particular, after observing a measured set \(S\) with outcome \(a\),
the remaining outcome probabilities are proportional to
\[
 \langle a,x|e^{-\beta\mathcal H}|a,x\rangle.
\]
For noncommuting \(\mathcal H\), they need not be the diagonal entries
of \(e^{-\beta\langle a|\mathcal H|a\rangle}\). Projecting the
Hamiltonian before exponentiation loses imaginary-time excursions
outside the measured subspace. A theorem stable under classical
Hamiltonian pinning therefore cannot be transferred to measured
quantum Gibbs laws by that substitution.

The quantum/classical comparison here concerns \emph{queries to a
classical string}. Complexity of preparing \(\rho_\beta\), evaluating
\(\operatorname{Tr}(O\rho_\beta)\), and sampling a new measurement
outcome are different tasks.
